\section{Chỗ khảo sát ấy đặt độ trung thực}
\label{sec:ch17-cho-dat-do-trung-thuc}

\index{độ trung thực!một ô bảng không có dụng cụ nào sau lưng|(}%
Mục trước cho thấy khảo sát xếp chuỗi suy luận vào nhóm không cần đối chiếu. Mục
này hỏi khảo sát ấy làm gì với chính chữ độ trung thực, vì nó có một mục dành
riêng cho việc đánh giá lời giải thích và trong mục ấy chữ này có mặt.

\index{độ trung thực!phép chia định tính với định lượng}%
Khảo sát chia việc đánh giá làm hai loại, định tính và định
lượng~\autocite{p16llmxai}. Bên định tính có hai tiêu chí, mức dễ hiểu với người
đọc và mức điều khiển được của lời giải thích. Bên định lượng cũng có hai, và
đó là \tn{độ trung thực}{faithfulness} với
\tn{tính hợp lý}{plausibility}. Phép chia này chính là phép chia mà
mục~\ref{sec:ch01-do-trung-thuc-va-tinh-hop-ly} dựng, nên tới đây thì khảo sát
đang đứng đúng chỗ.

\index{độ trung thực!ba điều kiện của khảo sát}%
Định nghĩa mà khảo sát đưa ra có ba điều kiện, và ba điều kiện ấy không cùng
nói về một quan hệ. Bài viết rằng một lời giải thích trung thực phải khớp với ít
nhất ba thứ: logic bên trong và quan hệ nhân quả mà mô hình xác lập; cách mô
hình được thiết kế và dữ liệu mô hình được huấn luyện trên; và việc quyết định
của mô hình có tái lập được bằng những đặc trưng mà lời giải thích nêu ra hay
không~\autocite{p16llmxai}.

Đọc ba điều kiện ấy theo bảng~\ref{tab:ch01-bon-quan-he} thì chúng rơi vào ba
chỗ khác nhau. Điều kiện thứ nhất là dòng độ trung thực, đúng
định nghĩa~\ref{def:ch01-do-trung-thuc}. Điều kiện thứ ba là một phép thử tái
lập, tức là hình dạng của họ sufficiency ở
mục~\ref{sec:ch08-ho-xoa-dan}, nên nó là một điều kiện cần chứ không phải một
định nghĩa. Điều kiện thứ hai thì không rơi vào dòng nào cả: khớp với cách mô
hình được thiết kế và với dữ liệu huấn luyện là một quan hệ giữa lời giải thích
và bối cảnh huấn luyện, và bảng~\ref{tab:ch01-bon-quan-he} không có dòng cho nó,
vì không có suy diễn nào về hành vi của mô hình đi ra từ nó. Một lời giải thích
hoàn toàn khớp với dữ liệu huấn luyện vẫn có thể mô tả sai cái mô hình đã dùng ở
một đầu vào cụ thể.

\index{độ trung thực!ô định lượng không có gì định lượng}%
Chỗ đáng ghi nhất thì nằm trong bảng 2 của khảo sát, trang 8. Bảng ấy có bốn
dòng, một dòng cho mỗi tiêu chí, với bốn cột: loại, tên tiêu chí, mô tả, và ví
dụ. Dòng độ trung thực nằm bên \enquote{Quantitative}, mô tả của nó là mức chính
xác trong việc biểu diễn quá trình ra quyết định bên trong mô hình, và ô ví dụ
của nó ghi một hệ chẩn đoán y khoa cho thấy những đặc trưng của bệnh nhân và các
trọng số của chúng~\autocite{p16llmxai}. Ô ví dụ của một tiêu chí định lượng mô
tả một bức tranh về một hệ thống, không phải một con số.

Đó không phải chuyện của riêng ô bảng. Rà cả 18 trang, khảo sát không đưa ra
công thức nào, ngưỡng nào, thủ tục nào hay điểm số nào để tính đại lượng mà bảng
2 gọi là độ trung thực. Nó cũng không ở đâu nói rằng chưa có cách tính. Chữ
\enquote{unfaithful} không xuất hiện lần nào trong toàn bài, nên trường hợp tiêu
chí ấy không đạt cũng không có tên. Và mục hạn chế của khảo sát, mục 7, liệt kê
năm khó khăn của việc giải thích bằng mô hình ngôn ngữ, gồm dữ liệu nhạy cảm,
chuẩn mực xã hội khác nhau, dữ liệu nhiều nguồn, độ phức tạp của mô hình, và
thiên lệch trong mô hình ngôn ngữ; không khó khăn nào trong năm khó khăn ấy là
việc không ai kiểm được lời giải thích có trung thực hay không.

\index{độ trung thực!nguồn chuẩn bị xếp sang dòng khác}%
Còn một chi tiết trong cùng bảng ấy và sách nêu nó hẹp đúng bằng cái nó là. Ô
nguồn của dòng tính hợp lý dẫn hai số hiệu, và một trong hai trỏ về bài của
Jacovi và Goldberg, nhan đề \emph{Towards faithfully interpretable NLP systems:
How should we define and evaluate faithfulness?}~\autocite{p16llmxai}. Đó là bài
mà mục~\ref{sec:ch09-dinh-nghia-lai} gọi là nguồn của định nghĩa được dùng rộng
rãi nhất về độ trung thực, và là bài mà định nghĩa~\ref{def:ch01-do-trung-thuc}
truy về. Trong bảng, nó đứng ở dòng bên cạnh.

Chỗ đặt ấy không phải một chỗ lệch riêng của cái bảng. Phần thân dẫn chính bài
ấy hai lần, và hai lần theo hai cách khác nhau. Lần thứ nhất, trang 6, dẫn đúng:
ở đó nó được dẫn cho ý rằng một lời giải thích trông dễ hiểu với người vẫn có thể
không đúng với logic thật của mô hình, tức là chính phép phân biệt giữa hai tiêu
chí. Lần thứ hai, trang 7, nằm trong đúng đoạn định nghĩa tính hợp lý, làm nguồn
cho ví dụ về một mô hình dự báo khí hậu đưa ra lời giải thích dựa trên các nguyên
lý khí tượng đã được thiết lập~\autocite{p16llmxai}. Nghĩa là nhãn dòng trong
bảng không mâu thuẫn với phần thân; nó theo đúng cách dùng thứ hai. Điều sách nói
vì thế hẹp hơn một lời buộc tội và cứng hơn một nhận xét về cái bảng: trong cả
bài, nguồn kinh điển của định nghĩa độ trung thực chỉ xuất hiện ở phía tính hợp
lý, trừ đúng một câu vẽ ranh giới giữa hai phía.

\index{độ trung thực!những con số duy nhất của khảo sát}%
Khảo sát có in số, và những con số ấy nằm ở một mục khác. Mục 8 so sánh các điểm
số saliency của sáu mô hình ngôn ngữ dưới hai độ dài ngữ cảnh, và giải thích rằng
kỹ thuật saliency giúp nhận ra những phần của một lời giải thích được sinh ra,
chẳng hạn từ hoặc cụm từ, khớp nhất với ground truth~\autocite{p16llmxai}.
Bảng~\ref{tab:ch17-saliency} chép lại mười hai con số ấy từ hình 4 của bài.

Bảng ấy nói ba điều. Thứ nhất, đại lượng được đo là mức khớp giữa những chỗ được
làm nổi bật trong một lời giải thích và một văn bản tham chiếu, nên theo
bảng~\ref{tab:ch01-bon-quan-he} nó là một đại lượng bên tính hợp lý; đó là cách
đọc của sách và bài không nói vậy. Thứ hai, phần thân gọi mô hình ở dòng cuối là
OLMO (1B) trong khi hình 4 của chính bài dán nhãn dòng ấy là Olmo (7B), và hai
con số thì khớp~\autocite{p16llmxai}. Thứ ba, bài đọc bảng này thành một câu về
quy mô, rằng khoảng cách giữa các mô hình cho thấy tầm quan trọng của kích thước,
kiến trúc và cách huấn luyện, và xếp Vicuna với Mistral vào cùng một nhóm điểm
thấp; nhưng Mistral ở 7 tỉ tham số đạt 0.68, cao hơn Vicuna ở 13 tỉ với 0.57 và
chỉ kém Cohere ở 54 tỉ với 0.71 đúng 0.03. Thứ hạng không đơn điệu theo quy mô,
và chỗ đảo nằm đúng giữa hai mô hình mà bài xếp chung nhóm.

\begin{table}[htbp]
  \centering
  \caption{Mười hai điểm số saliency của hình 4 trong khảo sát về mô hình ngôn
  ngữ cho XAI, chép theo đúng nhãn mà hình đặt cho từng dòng. Phần thân của bài
  gọi mô hình ở dòng cuối là OLMO (1B). Thứ hạng theo cột ngữ cảnh dài không đơn
  điệu theo số tham số: Mistral ở 7 tỉ đứng trên Vicuna ở 13 tỉ.}
  \label{tab:ch17-saliency}
  \begin{tabular}{@{}lrr@{}}
    \toprule
    Mô hình, theo nhãn của hình & Ngữ cảnh dài & Ngữ cảnh ngắn \\
    \midrule
    ChatGPT (175B)  & 0.84 & 0.56 \\
    Llamav2 (70B)   & 0.75 & 0.53 \\
    Cohere (54B)    & 0.71 & 0.56 \\
    Vicuna (13B)    & 0.57 & 0.51 \\
    Mistral (7B)    & 0.68 & 0.50 \\
    Olmo (7B)       & 0.45 & 0.29 \\
    \bottomrule
  \end{tabular}
\end{table}

\index{độ trung thực!một ô bảng không có dụng cụ nào sau lưng|)}%
Gộp lại thì khảo sát này đứng ở đúng chỗ mà bảy chương của Phần~IV mô tả, chỉ
khác là nó đứng đó cho một họ lời giải thích khác. Nó gọi tên tiêu chí, đặt tiêu
chí ấy vào cột định lượng, không cấp dụng cụ nào để tính, và những con số duy
nhất nó in ra thì đo một quan hệ khác. Điều mục sau đọc là một khảo sát nhìn thấy
chỗ hổng ấy và gọi tên ba họ dụng cụ cho nó.
